\subsection{牛顿定律应用的一般次序}

复杂装置仍由若干简单物体和约束组成。第一步确定哪些物体具有相同大小的加速度，哪些方向
必须分别处理；第二步对每个研究对象画受力图；第三步选轴并列出分量方程；最后才代入运动学
关系。在忽略绳质量、滑轮转动惯量和轮轴摩擦，且绳上没有其他切向分布力的理想模型中，
同一根绳中的张力才可视为处处相等。绳长不变则给出绳各段位移、速度和加速度之间的几何约束；
有动滑轮时须把所有可变绳段计入，不能一概认为两端加速度大小相等。
绳或滑轮有质量时，张力通常不再相等，但只要绳仍不可伸长，绳长约束仍成立。

摩擦力尤其不能先写成 $\mu N$。静摩擦力在零到最大值之间自适应，只有由平衡或动力学方程
算出的需要值满足 $|f_s|\leq\mu_sN$，原先假定的相对静止才可能成立；若超过极限，物体发生
滑动，此时才用 $f_k=\mu_kN$，方向与接触面的相对滑动速度相反。在叠放物块、传送带和斜面
问题中，应先假定一种运动状态，算出所需摩擦，再用不等式检验，必要时分段讨论。

圆周运动也不对应一种新的力。“向心力”只是合力的法向分量，大小为
\begin{equation}
 F_n=m\frac{v^2}{R}=m\omega^2R.
\end{equation}
重力、支持力、张力或摩擦力中的某些分量共同提供它。竖直圆周运动中速率随高度改变，最低点
与最高点的受力方程不能混用；轻绳模型还要检查张力不能为负，轨道约束则要检查支持力不能
指向轨道不可能施力的一侧。这些条件常常决定物体能否完成整周运动。

万有引力问题把上述思路推广到距离变化的中心力。球对称天体外部可等效为质量集中在球心，
$F=GMm/r^2$；近地面小范围内才可把 $g$ 当作常量。月球的向心加速度与地表重力的比较说明，
落体与天体运动可以由同一种相互作用解释。计算以后还应检查量纲、方向和极限情形，避免只得到
形式正确却物理上不可能的根。
